\section{Einführung und Ziele}

Dieses Kapitel beschreibt die wesentlichen Anforderungen und Treiber, die bei der Architekturentscheidung berücksichtigt wurden: Aufgabenstellung, Geschäftskontext, Qualitätsziele und relevante Stakeholder.

\subsection{Aufgabenstellung}

Die SkillBridge AG ist ein mittelgrosses Schweizer Unternehmen mit Sitz in Biel/Bienne~(BE), das sich auf die Digitalisierung von Weiterbildungsangeboten für Berufsverbände, Bildungsinstitutionen und private Kursanbieter spezialisiert. Das Kernprodukt ist \textbf{EduPlanner} -- eine webbasierte Weiterbildungsplattform, die Kursanbieter bei der vollständigen Verwaltung ihres Kursgeschäfts unterstützt.

EduPlanner bildet den gesamten Lebenszyklus eines Kursangebots ab -- von der Erstellung und Veröffentlichung über die Planung konkreter Durchführungen bis hin zur Einschreibung von Teilnehmern und der automatischen Kommunikation relevanter Ereignisse.

\subsubsection{Geschäftskontext und Marktposition}

EduPlanner positioniert sich als spezialisierte SaaS-Lösung für mittelständische Kursanbieter
mit 10 bis 500 Mitarbeitenden. Im Wettbewerbsumfeld konkurriert die Plattform primär mit drei
Alternativen: der manuellen Verwaltung über Excel oder Google Sheets, kostenintensiven
Legacy-Systemen sowie generischen Learning Management Systems (LMS), denen dedizierte
Kursplanungs-Funktionalitäten fehlen.


Gegenüber diesen Alternativen differenziert sich EduPlanner durch vier Merkmale, die
gleichzeitig zentrale Treiber der Architekturentscheidungen darstellen:

\begin{itemize}[leftmargin=1.5em]
  \item \textbf{Benutzerfreundlichkeit:} Eine intuitive Oberfläche für nicht-technisches Personal
  erfordert klare UI-Richtlinien – konkret: maximal 3~Klicks pro Kernfunktion und
  Ladezeiten unter 500\,ms.

  \item \textbf{Schnelle Implementierung:} Das Ziel, innerhalb von 2~Wochen produktiv zu gehen,
  macht eine vollständig automatisierte CI/CD-Pipeline zur Voraussetzung.

  \item \textbf{Schweizer Datenschutz:} Die Einhaltung von DSGVO und DSG erfordert einen
  Privacy-by-Design-Ansatz sowie eine Datenspeicherung ausschließlich in der EU
  bzw.\ der Schweiz.

  \item \textbf{API-First:} Um die Integration in bestehende Systemlandschaften zu ermöglichen,
  sind eine vollständige OpenAPI-Dokumentation und eine klare Versionierungsstrategie
  verbindlich.
\end{itemize}


\subsubsection{Erwartete Nutzungszahlen (Jahr~1)}
\begin{itemize}[leftmargin=1.5em]
  \item 5--20 Kursanbieter (Tenants)
  \item 10--50 aktive Kurse im Katalog pro Anbieter
  \item 50--500 Einschreibungen pro Monat
  \item Bis zu 100 gleichzeitige Nutzer zu Spitzenzeiten
\end{itemize}

\subsubsection{Kernfunktionen des Systems}

\begin{tabularx}{\textwidth}{L{3cm} X}
  \rowcolor{tableheader}
  \color{white}\textbf{Funktion} & \color{white}\textbf{Beschreibung} \\
  Kurskatalog & Kurse anlegen, beschreiben, veröffentlichen und archivieren; Trainer zuordnen \\
  \rowcolor{tablegrey}
  Kursplanung & Konkrete Durchführungen (Termine, Orte, Kapazitäten) planen und verwalten \\
  Einschreibung & Teilnehmer einschreiben, Wartelistenverwaltung, Stornierungen \\
  \rowcolor{tablegrey}
  Benachrichtigungen & Automatische E-Mails bei Einschreibung, Wartelistenaufnahme und Stornierung \\
  Sicherheit & Rollenbasierte Zugriffskontrolle für alle Benutzergruppen \\
  \rowcolor{tablegrey}
  Reporting & Dashboard mit Auslastungsquoten, Teilnehmerstatistiken und Kursübersicht \\
\end{tabularx}

\subsubsection{Go-Live-Kriterien}

\textbf{MVP (Minimum Viable Product) -- Monat 1--6:}
\begin{itemize}[leftmargin=1.5em]
  \item \textbf{Kurskatalog-Verwaltung:} Kurse können erstellt, bearbeitet,
        veröffentlicht und gelöscht werden (vollständige CRUD-Operationen)
  \item \textbf{Kursplanung:} Termine, Durchführungsorte und Teilnehmerkapazitäten
        können pro Kurs erfasst und bearbeitet werden
  \item \textbf{Einschreibung:} Teilnehmer können sich einschreiben; bei voller
        Kapazität erfolgt automatische Aufnahme in die Warteliste mit
        Nachrücklogik
  \item \textbf{E-Mail-Benachrichtigungen:} Automatischer Versand bei Einschreibung,
        Wartelistenaufnahme und Stornierung
  \item \textbf{Zugriffskontrolle:} Vier Rollen (Admin, Kursanbieter, Trainer,
        Teilnehmer) mit klar getrennten Berechtigungen sind vollständig
        implementiert und getestet
  \item \textbf{Verfügbarkeit:} System erreicht nachweislich 99,5\,\% Uptime
        zu Kernzeiten im Staging-Betrieb (vgl.\ Q2)
  \item \textbf{Überbuchungsschutz:} Kein einziger Fall einer Überbuchung in
        den Integrationstests nachgewiesen (vgl.\ Q1)
\end{itemize}

\textbf{v1.0 (Vollständiger Release) -- Monat 7--12:}
\begin{itemize}[leftmargin=1.5em]
  \item \textbf{Erweiterte Planung:} Unterstützung mehrerer Durchführungen pro Kurs
  (Multi-Execution) inkl.\ Konfliktprüfung bei Trainer- und Raumzuweisung
  \item \textbf{Reporting:} Auslastungsberichte, Einschreibungsstatistiken und
  Export-Funktion (CSV/PDF) für Kursanbieter
  \item \textbf{API-Dokumentation:} Vollständige OpenAPI~3.0-Spezifikation für alle
  öffentlichen Endpunkte, inkl.\ Beispielaufrufen und Fehlercode-Beschreibungen
  \item \textbf{Performance-Verifikation:} Lasttest mit 1\,000 gleichzeitigen Nutzern
  bestanden; Ladezeiten unter 500\,ms (Q6, Q11)
  \item \textbf{Security-Audit:} Penetrationstest gemäss OWASP Top~10 durchgeführt
  und alle kritischen Befunde behoben (Q5, Q7)
  \item \textbf{Datenschutz-Audit:} Nachweis der DSGVO/DSG-Konformität durch
  externen Audit bestätigt; Datenspeicherung ausschliesslich EU/Schweiz (Q12)
\end{itemize}

% ── 1.2 Qualitätsziele ────────────────────────────────────────
\subsection{Qualitätsziele}

\begin{tabularx}{\textwidth}{C{0.5cm} L{4.5cm} C{1.8cm} X}
  \rowcolor{tableheader}
  \color{white}\textbf{\#} & \color{white}\textbf{ISO-25010-Attribut\footnotemark} & \color{white}\textbf{Priorität} & \color{white}\textbf{Motivation} \\
  1 & Functional Correctness & Kritisch & Einschreibungen und Kapazitätsprüfungen müssen stets konsistent sein. Eine Überbuchung ist geschäftskritisch und nicht tolerierbar. \\
  \rowcolor{tablegrey}
  2 & Reliability -- Availability & Hoch & Das System muss zu Kernzeiten (Mo--Fr, 07:00--22:00~Uhr) mit $\geq$\,99,5\,\% Verfügbarkeit betrieben werden. \\
  3 & Maintainability & Hoch & Fachliche Domänen müssen unabhängig voneinander entwickelt, getestet und deployt werden können. \\
  \rowcolor{tablegrey}
  4 & Security & Hoch & RBAC muss sicherstellen, dass Nutzer ausschliesslich auf freigegebene Funktionen und Daten zugreifen. \\
  5 & Usability & Mittel & Kernprozesse in $\leq$\,3~Klicks; Ladezeiten $<$\,500\,ms. \\
  \rowcolor{tablegrey}
  6 & Performance Efficiency & Mittel & P95-Ladezeit unter 500\,ms bei bis zu 1\,000 gleichzeitigen Nutzern (vgl.\ Q6, Q11). \\
\end{tabularx}
\footnotetext{ISO/IEC~25010 ist ein internationaler Standard, der Softwarequalität
in acht übergeordnete Merkmale gliedert (u.\,a.\ \textit{Performance Efficiency},
  \textit{Security}, \textit{Maintainability}, \textit{Reliability}). Jedes Merkmal
  ist wiederum in konkrete Teilmerkmale unterteilt, die als messbare Qualitätsziele
  in der Architektur verankert werden können.}




% ── 1.3 Stakeholder ───────────────────────────────────────────
\subsection{Stakeholder}

\begin{tabularx}{\textwidth}{L{2.2cm} L{3.5cm} X}
  \rowcolor{tableheader}
  \color{white}\textbf{Rolle} & \color{white}\textbf{Name / Funktion} & \color{white}\textbf{Erwartungen / Interessen} \\
  Auftraggeber & Sandra Meier, CEO &
  Erwartet messbaren ROI durch Reduktion manueller Verwaltungsaufwände.
  Das System soll mit wachsender Kundenzahl skalieren, ohne proportional
  mehr Personal zu erfordern. Entscheidet über Budget und Go/No-Go. \\
  \rowcolor{tablegrey}
  Product Owner & Thomas Zbinden, Leiter Produktentwicklung &
  Verantwortet die termingerechte Lieferung aller vereinbarten Features.
  Priorisiert den Backlog, nimmt Sprints ab und stellt sicher, dass
  fachliche Anforderungen korrekt umgesetzt werden. \\
  Fachbereich & Kursplanungs-Team &
  Nutzt das System täglich für die operative Planung. Erwartet eine
  intuitive Oberfläche ohne Schulungsaufwand, eine übersichtliche
  Kalenderansicht sowie schnelle Reaktionszeiten bei der Dateneingabe. \\
  \rowcolor{tablegrey}
  Endbenutzer & Teilnehmer &
  Benötigen einen reibungslosen Self-Service: Kurssuche, Einschreibung
  und Stornierung in wenigen Klicks. Erwarten zuverlässige
  Bestätigungs- und Erinnerungsbenachrichtigungen per E-Mail. \\
  Endbenutzer & Trainer (Kursleiter) &
  Benötigen eine persönliche Übersicht ihrer zugewiesenen Durchführungen,
  Termine und Teilnehmerlisten. Erwarten rechtzeitige Benachrichtigung
  bei Änderungen oder Stornierungen. \\
  \rowcolor{tablegrey}
  Betrieb & Patrick Hari, IT-Leiter &
  Verantwortet den stabilen Betrieb der Plattform. Erwartet einfaches
  Deployment via CI/CD, aussagekräftiges Monitoring und Alerting sowie
  klar dokumentierte Konfigurationsparameter ohne Eingriffe ins
  Quellcode. \\
  Entwicklung & Dev-Team &
  Benötigt eine klar strukturierte, gut dokumentierte Architektur mit
  hoher Testabdeckung. Erwartet automatisierte CI/CD-Pipelines, saubere
  Schnittstellendefinitionen zwischen den Services und geringe
  Kopplungen zwischen den Bounded Contexts. \\
  \rowcolor{tablegrey}
  Datenschutz & Andrea Senn, CISO &
  Stellt die Einhaltung von DSGVO und Schweizer DSG sicher. Erwartet
  Privacy-by-Design, eine revisionssichere Protokollierung von
  Datenzugriffen sowie eine starke Authentifizierung und
  rollenbasierte Zugriffskontrolle für alle personenbezogenen Daten. \\
\end{tabularx}